\documentclass{article}
\usepackage{pathfinder-readable}

\title{What Can Be Certified in a Room-Guided Door Graph?}

\begin{document}
\maketitle

\begin{abstract}
This note connects a survey of learning-augmented algorithms with a robot that builds a scene graph from room
and door concepts. The thread asked whether the robot's claims about which rooms connect could be checked when
its room model also guides what it observes. The peers proved stability for one door-detection threshold under
a stated wall-error bound, and a sensing lower bound in a separate, simplified audit model. Neither result
certifies the robot's complete graph or bounds its physical sensing cost. The verifier paused the thread
because a usable claim needs a LiDAR visibility and error model, followed by experiments that test whether
wall regions can be resolved at reasonable cost.
\end{abstract}

\section{Introduction}

Q surveys algorithms that use fallible predictions with formal performance guarantees \cite{Q}. It explains
why component guarantees need separate assumptions when predictions change later observations, instances or
costs. Transferring an error bound needs a metric and a bound on sensitivity. Its charge for prediction calls
does not price robot travel or scans.

P describes a robot that builds an indoor scene graph from room and door concepts, without a prior global
metric map \cite{P}. Once it accepts a room, its walls guide the door search. A detected door can prompt a
visit and then a crossing. Simplified tests recover the correct room connections, but false rooms can be
accepted and later revision is not implemented.

The connection pursued was the room-to-door dependency. Can a slightly wrong room support a door decision? Can
an unseen wall support a claim that no door is there? These questions join Q's concerns about error and
missing feedback to P's detector.

\section{What was done}

The peers first traced how P's estimated walls filter LiDAR points for door detection. They compared this hard
filter with Q's conditions for transferring an error guarantee. A small wall change can flip a point near the
threshold. Ada derived a stability condition for the first filter with a fixed scan.

Emmy then modelled hidden openings as wall regions checked one by one. Her coupling argument measured the cost
of claiming exact topology even with a correct no-door prediction. Ada checked it and noted that a real LiDAR
view can cover several regions.

The peers finally compared evidence in P's graph. A door becomes nominal after a visit, before a crossing. A
crossing witnesses a passage, but naming its rooms needs correct recognition and association. P also describes
geometric loop closure through an uncrossed door. The peers proposed recording candidates, matches, crossings
and unresolved regions separately. P does not supply this certificate.

\section{Results and their scope}

\begin{proposition}
Fix a finite scan \(S\), an estimated wall set \(\widehat W\) and a true wall set \(W^\star\) in the same
coordinates. Let \(\tau\) be P's wall-distance threshold, \(d_H\) the Hausdorff distance between wall sets,
and \(\varepsilon\) a certified upper bound on their error. Define the smallest observed threshold margin by
\[
m=\min_{p\in S}\left|\operatorname{dist}(p,\widehat W)-\tau\right|.
\]
If \(d_H(W^\star,\widehat W)\leq\varepsilon<m\), every point in \(S\) makes the same pass or fail decision at
the first wall-distance filter under either wall set.
\end{proposition}

Here \(\operatorname{dist}(p,W)\) is the distance from scan point \(p\) to wall set \(W\). It changes by at
most the Hausdorff distance when the wall set changes. Thus a margin larger than the certified error prevents
a threshold crossing. Ada's derivation in \texttt{ada/margin-gated-scene-graphs.md} applies this to P's
\(0.15\,\mathrm m\) filter. At the threshold, tiny wall changes can reverse the decision. Hence this binary
gate has no global finite Lipschitz bound under wall distance, though Q allows other error-dependent
guarantees.

\begin{proposition}
In an audit-only model with \(n\) hidden wall regions, suppose each region has either a door or no door.
Checking one region costs one unit and reveals only that region's state. Let \(\delta\), with
\(0\leq\delta<1/2\), be the permitted failure probability. If an algorithm must report every region correctly
with probability at least \(1-\delta\) for every possible state under the same prediction, then, on the state
with no doors, its expected number of distinct checks is at least \(n(1-2\delta)\).
\end{proposition}

For any region \(j\), compare a no-door state with a state whose only door is in \(j\). Give the algorithm the
same prediction and random choices in both runs. Until it checks \(j\), the runs look identical. If it omits
\(j\) and correctly reports no doors in the first run, it must be wrong in the second. The allowed failure
probability in each run therefore bounds the chance of omitting \(j\) by \(2\delta\); summing over regions
gives the result. Emmy's proof is in \texttt{emmy/verification\_cost.md}, and Ada checked it independently.
The bound concerns independent unit-cost checks, not P's LiDAR views, travel, or sensing policy.

The results address different gaps. A margin supports a decision on measured points; an absence claim needs
coverage and a sensor model. Neither proves P's graph correct. A crossing witnesses passage, while a named
room pair also needs reliable identities and association.

\section{Objections and unresolved questions}

P's other door tests include derivative peaks, door widths, candidate pairing and association. The
first-filter margin gives no bound for these decisions or for asynchronous graph updates. P does not provide
the deterministic wall-error bound assumed by the proposition. Its correct topology in simplified tests is
evidence for those tests, not a general certificate. Likewise, the audit bound cannot be read as a cost bound
for a LiDAR robot: one view may check many regions, and Q's charge for predictor calls does not price travel
or sensing.

An unresolved label would avoid false claims if applied to every edge, but it would tell the robot little. A
useful result needs a stated resolution target and an observation-cost measure. It also needs a rule for
revising door claims when an accepted room changes. P proposes later fitness monitoring and alternative
hypotheses, but does not implement that revision in the reported system \cite{P}.

The peers made a bounded check of related abstracts. \emph{Active Semantic Perception} and \emph{SCOUT}
already connect uncertain scene graphs with active traversal \cite{active-semantic-perception,scout}.
\emph{Estimating Map Completeness in Robot Exploration} addresses a learned estimate of remaining map
coverage, while \emph{Landmark-based Distributed Topological Mapping and Navigation} uses designed landmark
coverage conditions \cite{map-completeness,landmark-mapping}. Those checks show adjacent work; they do not
establish novelty or exhaust the literature.

\section{Why the thread stopped}

The verifier accepted the two elementary arguments and their limits, but paused because neither yields a
connectivity certificate or sensing-cost bound for P's robot. The smallest next step is to specify visibility
and wall-error assumptions for one room, then test P's room-and-door scenario with wall perturbations near the
threshold and occluded openings. The test should record asserted, absent and unresolved edges, wrong room-pair
connections, and the sensing and travel needed to resolve wall regions. A useful certificate for that
restricted setting would restart the larger question; the present record does not establish one.

\bibliographystyle{plainurl}
\bibliography{references}
\end{document}
